Research on misleading visualizations connects two
lines.
The first studies how charts mislead human readers:
deception relies less on outright distortion than on
faithful drawings with selective context and framing
\citep{lisnic2023misleading, nadib2026guardrail}.
The second measures the phenomenon in models:
benchmarks evaluate multimodal LLMs on misleading chart
QA \citep{chen2025unmasking, bharti2024chartom,
alexander2024models}, truncation, aspect-ratio, and
encoding manipulations are quantified with mitigations
\citep{mahbub2025perils, mahbub2026chart}, and recent
work asks whether agents recognize misleadingness
unprompted \citep{panda2026making, blasilli2026true,
seto2026llms}.
Method-side responses include tool-augmented chart
reasoning \citep{kaur2025chartagent}, structured
extraction \citep{das2025chartsofthought},
misleading-chart detection \citep{tonglet2025this}, and
agentic dual-path frameworks \citep{zhang2026navigating}.
\citet{tonglet2025protecting} provide the closest
model-level defense through tabular answering and
re-plotting; these repair a single question--answer
round, while \ours{} targets the interpretation rules
that recur across an executing agent's rounds and
maintains their validated revisions in runtime state.

A second line concerns agent-environment integrity:
injection benchmarks show that tool and web content can
steer agents \citep{zhan2024injecagent}, security
benchmarks measure targeted hijacking
\citep{evtimov2025wasp}, injection reaches perceived
representations such as HTML accessibility trees
\citep{johnson2025manipulating}, and red-teaming and
runtime frameworks generate and intercept such attacks
at scale \citep{wang2025agentvigil, zhao2026clawguard}.
Interface-level defenses separate instructions from data
\citep{chen2024struq} or train against injected content
\citep{chen2024secalign}.
\citet{wu2024dissecting} dissect the robustness of
multimodal agents as compound systems and show that
reflection and tree search, the components commonly
added for self-correction, open new vulnerabilities when
the evaluator itself is exposed to attacked content.
Our defense differs in what it protects: not the agent's
instruction hierarchy but the interpretation layer
between visual observation and task judgment, where a
misleading chart exerts its effect without any injected
instruction.

The third line concerns persistent state and
correction: verbal reinforcement carries textual
feedback across attempts \citep{shinn2023reflexion},
self-consistency aggregates reasoning paths
\citep{wang2022selfconsistency}, and intrinsic
self-correction without external feedback degrades
reasoning \citep{huang2023large, li2025decomposing};
agent surveys organize the state and memory design
space \citep{wang2024survey, hu2025agents}, and
realistic visual web tasks remain hard enough that
single-decision correctness cannot be taken for
granted \citep{koh2024visualwebarena}.
\ours{} borrows the state-carrying idea but changes the
carried content: the state records interpretation rules
with evidential support, applicability scope, and
revision history, so later steps inherit what was
verified about the chart rather than a verdict about
the answer.
